\documentclass[11pt,twoside]{article}
% Shared formatting for all notes in this repository.
% Adapted from the UC Berkeley EECS CS267 lecture-note template:
% https://www.cs.cmu.edu/~ggordon/10725-F12/template.tex
% Personal-note adaptation: boxed title/author/date; ordinary numbering;
% existing natbib citations and mathematical environments are retained.
% Load after \documentclass; compile from the repository root.
\usepackage[T1]{fontenc}
\usepackage[utf8]{inputenc}
\usepackage{amsmath,amsfonts,amssymb,mathrsfs,amsthm}
\usepackage{graphicx,wrapfig}
\usepackage[inline,shortlabels]{enumitem}
\usepackage{centernot}
\usepackage[numbers]{natbib}
\usepackage{adjustbox,booktabs,arydshln}
\usepackage{caption,subcaption}
\usepackage{xcolor,tikz,pifont}
\usetikzlibrary{decorations.pathreplacing,angles,quotes,shapes.geometric,
  arrows.meta,positioning,calc,fit,backgrounds,shadows}
\usepackage[ruled,vlined]{algorithm2e}
\usepackage{xr}
\usepackage{hyperref}
\hypersetup{colorlinks=true,linkcolor=blue,citecolor=blue,urlcolor=blue}

% Berkeley EECS lecture-note page layout and paragraph spacing.
\setlength{\oddsidemargin}{0.25in}
\setlength{\evensidemargin}{-0.25in}
\setlength{\topmargin}{-0.6in}
\setlength{\textwidth}{6.5in}
\setlength{\textheight}{8.5in}
\setlength{\headsep}{0.75in}
\setlength{\parindent}{0in}
\setlength{\parskip}{0.1in}

% A wrapping box accommodates long research-note titles.
\makeatletter
\renewcommand{\maketitle}{%
  \pagestyle{myheadings}%
  \thispagestyle{plain}%
  \begin{center}
    \fbox{%
      \begin{minipage}{\dimexpr\textwidth-2\fboxsep-2\fboxrule\relax}
        \vspace{2mm}
        \textbf{Research Notes}\hfill\@date\par
        \vspace{2mm}
        {\centering\Large\bfseries\@title\par}
        \vspace{2mm}
        {\itshape Author: \@author\par}
        \vspace{2mm}
      \end{minipage}%
    }
  \end{center}
  \markboth{Research Notes}{Research Notes}%
  {\footnotesize\itshape
    Style adapted from the UC Berkeley EECS lecture-note template.\par}
  \vspace{4mm}%
}
\makeatother

\newcommand{\bR}{{\boldsymbol R}}
\newcommand{\bX}{{\boldsymbol X}}
\newcommand{\bY}{{\boldsymbol Y}}
\newcommand{\by}{{\boldsymbol y}}
\newcommand{\bx}{{\boldsymbol x}} 
\newcommand{\bz}{{\boldsymbol z}}
\newcommand{\bZ}{{\boldsymbol Z}}
\newcommand{\bH}{{\boldsymbol H}}
\newcommand{\bb}{{\boldsymbol b}}
\newcommand{\be}{{\boldsymbol e}}
\newcommand{\bbf}{{\boldsymbol f}}
\newcommand{\bg}{{\boldsymbol g}}
\newcommand{\bm}{{\boldsymbol m}}
\newcommand{\bw}{{\boldsymbol w}}
\newcommand{\bW}{{\boldsymbol W}}
\newcommand{\bv}{{\boldsymbol v}}
\newcommand{\bt}{{\boldsymbol t}}
\newcommand{\bA}{{\boldsymbol A}}
\newcommand{\ba}{{\boldsymbol a}}
\newcommand{\bD}{{\boldsymbol D}}
\newcommand{\bE}{{\boldsymbol E}}
\newcommand{\bi}{{\boldsymbol i}}
\newcommand{\bc}{{\boldsymbol c}}
\newcommand{\bC}{{\boldsymbol C}}
\newcommand{\bk}{{\boldsymbol k}}
\newcommand{\bK}{{\boldsymbol K}}
\newcommand{\bP}{{\boldsymbol P}}
\newcommand{\mK}{\mathcal{K}}
\newcommand{\bV}{{\boldsymbol V}}
\newcommand{\bG}{{\boldsymbol G}}
\newcommand{\mG}{\mathcal{G}}
\newcommand{\mH}{\mathcal{H}}
\newcommand{\mE}{\mathbb{E}}
\newcommand{\mX}{\mathcal{X}}
\newcommand{\mM}{\mathcal{M}}
\newcommand{\bU}{{\boldsymbol U}}
\newcommand{\bu}{{\boldsymbol u}}
\newcommand{\bI}{{\boldsymbol I}}
\newcommand{\bJ}{{\boldsymbol J}}
\newcommand{\bs}{{\boldsymbol s}}

\newcommand{\bxi}{{\boldsymbol \xi}}
\newcommand{\balpha}{{\boldsymbol \alpha}}
\newcommand{\bbeta}{{\boldsymbol \beta}}
\newcommand{\bgamma}{{\boldsymbol \gamma}}
\newcommand{\btheta}{{\boldsymbol \theta}}
\newcommand{\bTheta}{{\boldsymbol \Theta}}
\newcommand{\bvartheta}{{\boldsymbol \vartheta}}
\newcommand{\bepsilon}{{\boldsymbol \epsilon}}
\newcommand{\bvarepsilon}{{\boldsymbol \varepsilon}}
\newcommand{\bomega}{{\boldsymbol \omega}}
\newcommand{\bmu}{{\boldsymbol \mu}}
\newcommand{\bSigma}{{\boldsymbol \Sigma}}

\newcommand{\corr}{\mbox{Corr}}
\newcommand{\var}{{\mbox{Var}}}
\newcommand{\Var}{{\mbox{Var}}}
\newcommand{\cov}{{\mbox{Cov}}}
\newcommand{\diag}{{\mbox{diag}}}
\newcommand{\bUmis}{{\boldsymbol U}_{{\rm mis}}}
\newcommand{\bYmis}{{\boldsymbol Y}_{{\rm mis}}}
\newcommand{\bYobs}{{\boldsymbol Y}_{{\rm obs}}}
\newcommand{\tbXmis}{{\tilde{\boldsymbol X}}_{{\rm  mis}}}
\newcommand{\tbxmis}{{\tilde{\boldsymbol x}}_{{\rm  mis}}}
\newcommand{\bXmis}{{\boldsymbol X}_{{\rm mis}}}
\newcommand{\bXobs}{{\boldsymbol X}_{{\rm obs}}}
\newcommand{\bZmis}{{\boldsymbol Z}_{{\rm mis}}}
\newcommand{\bZobs}{{\boldsymbol Z}_{{\rm obs}}}
\newcommand{\tXmis}{{\tilde{X}}_{{\rm  mis}}}
\newcommand{\txmis}{{\tilde{x}}_{{\rm  mis}}}
\newcommand{\Xmis}{X_{{\rm  mis}}}
\newcommand{\Xobs}{X_{{\rm  obs}}}
\newcommand{\bymis}{{\boldsymbol y}_{{\rm  mis}}}
\newcommand{\byobs}{{\boldsymbol y}_{{\rm  obs}}}
\newcommand{\ymis}{y_{{\rm  mis}}}
\newcommand{\yobs}{y_{{\rm  obs}}}

\newcommand{\JRSSB}{Journal of the Royal Statistical Society: Series B}
\newcommand{\indep}{\rotatebox[origin=c]{90}{$\models$}}
\newcommand{\red}{\textcolor{red}}
\newcommand{\blue}{\textcolor{blue}}
\newcommand{\notetodo}[1]{\textcolor{blue}{[TODO: #1]}}

% Preserve the existing independent, section-based theorem counters.
\newtheorem{theorem}{Theorem}[section]
\newtheorem{lemma}{Lemma}[section]
\newtheorem{corollary}{Corollary}[section]
\newtheorem{formula}{Formula}[section]
\newtheorem{proposition}{Proposition}[section]
\theoremstyle{definition}
\newtheorem{definition}{Definition}[section]
\newtheorem{example}{Example}[section]
\newtheorem{assumption}{Assumption}
\theoremstyle{remark}
\newtheorem{remark}{Remark}

% Journal abbreviations used by the existing bibliography.
\def\AJE{American Journal of Epidemiology}
\def\ANNALS{Annals of Statistics}
\def\BIOK{Biometrika}
\def\BMCS{Biometrics}
\def\STATMED{Statistics in Medicine}
\def\STATSCI{Statistical Science}
\def\JSPI{Journal of Statistical Planning \& Inference}
\def\JRSSB{Journal of the Royal Statistical Society, Series B}
\def\COMMS{Communications in Statistics, Series A}
\def\JQT{Journal of Quality Technology}
\def\TECH{Technometrics}
\def\JASA{Journal of the American Statistical Association}
\def\CSDA{Computational Statistics \& Data Analysis}
\def\JCGS{Journal of Computational and Graphical Statistics}
\def\JCB{Journal of Computational Biology}
\def\BIOINF{Bioinformatics}
\def\JAMA{Journal of the American Medical Association}
\def\JNUTR{Journal of Nutrition}
\def\LETTERS{Letters in Probability and Statistics}
\def\JABES{Journal of Agricultural, Biological and Environmental Statistics}
\def\SCAN{Scandinavian Journal of Statistics}
\def\PNAS{Proceedings of the National Academy of Sciences USA}
\def\JMLR{Journal of Machine Learning Research}
\def\JSCS{Journal of Statistical Computation and Simulation}


\title{STAT 695 Statistical Learning Theory}
\author{Wen-Hung Wang}
\date{\today}

\begin{document}
\maketitle

\section{Minimax Rates}
The minimax loss is best accuracy under the worst case, defined as
    \[
    R_n = \inf_{\hat{\theta}} \sup_{\theta \in \Theta} E_{\theta}[\ell(\theta, \hat{\theta}(X_1,...,X_n))]
    \]
We wish to find the order of this rate. Generally speaking, one
\begin{enumerate}
    \item Upper bounds $R_n$ by a good estimator, e.g., MLE.
    \item Lower bounds $R_n$ as in Le Cam's method (uniform prior trick, divergence inequalities, etc.).
    \item Matches the order of upper and lower bounds, e.g,
    \[
    R_n \leq C_1 r_n, \qquad R_n \geq C_2 r_n
    \]
    then $R_n \asymp r_n$.
\end{enumerate}

\section{Le Cam's Method}
Le Cam's approach gives a lower bound on $R_n$. Here I only summarize some tricks:
\begin{enumerate}
    \item First consider a \textit{binary parameter space} $\Theta = \{\theta_1, \theta_2\}$. Set $X$ from $P_{\theta_1}$ as null and $X$ from $P_{\theta_2}$ as alternative.
    \item Put a uniform prior on $\theta$ and obtain a lower bound containing the total variation distance of $P_{\theta_1}$ and $P_{\theta_2}$ for \textit{test risk} $P_{\theta}(\phi(X) \neq \theta)$ where $\phi$ is a binary test.
    \item Generalize this lower bound to $X_1,...,X_n$ and convert TV to KL or Hellinger. For iid data, KL and H are easier to handle .
    \item Define the \textit{estimation risk}, which is the main goal
    \[
    \inf_{\hat{\theta}}\sup_{\theta \in \{\theta_1, \theta_2\}}E_{\theta}[\|\hat{\theta}-\theta\|^2].
    \]
    Then lower bound this using Chebyshev's inequality and define a binary test $\phi$ to link estimation risk and test risk.
    \item Apply the inequality obtained in 3 to conclude the proof.
    \item Use second-order approximation for KL and H to find the exact order, under regularity of $P_{\theta}$.
    \item Lastly, upper bound $\inf_{\hat{\theta}}\sup_{\theta \in \{\theta_1, \theta_2\}}E_{\theta}[\|\hat{\theta}-\theta\|^2]$ by considering taking supremum over a continuous $\Theta$ to conclude the proof. 
\end{enumerate}
Therefore, Le Cam's method gives
\[
R_n = \inf_{\hat{\theta}}\sup_{\theta \in \Theta} E_{\theta}[\|\hat{\theta}(X_1,...,X_n)-\theta \|^2] \geq \underbrace{\max_{\theta_1,\theta_2}\|\theta_1-\theta_2\|^2}_{\text{param separation}} \underbrace{D(n, P_{\theta_1}, P_{\theta_2})}_{\text{stats indistingushability}},
\]
which is of order $1/n$ under regularity conditions for $P_{\theta}$. The problem for this method is that the inequality used in 7 is loose
\[
\sup_{\theta \in \Theta}...\geq \sup_{\theta \in \{\theta_1, \theta_2\}} ...
\]
because it does not take dimension of $\Theta$ into account.

\section{Assouad's method}

\section{Fano's method}

\section{Summary of Minimax Lower Bound Methods}
The common strategy behind Le Cam's, Assouad's, and Fano's methods is to lower-bound the full minimax risk
\[
R_n
=
\inf_{\hat{\theta}}
\sup_{\theta\in\Theta}
E_{\theta}
L(\hat{\theta},\theta)
\]
by restricting attention to a carefully chosen hard subset
\[
\Theta_0\subseteq\Theta.
\]
Since
\[
\sup_{\theta\in\Theta}
E_{\theta}L(\hat{\theta},\theta)
\ge
\sup_{\theta\in\Theta_0}
E_{\theta}L(\hat{\theta},\theta),
\]
any lower bound proved on $\Theta_0$ is automatically a valid lower bound for the original problem.

The estimation problem is then converted into a testing or classification problem. Roughly,
\[
\boxed{
\text{if an estimator is very accurate}
\;\Longrightarrow\;
\text{it should identify which hypothesis generated the data}.
}
\]
Therefore, if the corresponding testing problem is hard, estimation must also be hard.

The three methods differ mainly in the structure of the hard subset:
\[
\boxed{
\begin{array}{c|c}
\text{Le Cam} & \text{two points} \\[2pt]
\text{Assouad} & \{-1,+1\}^m \text{ hypercube} \\[2pt]
\text{Fano} & M \text{ well-separated points / packing}
\end{array}
}
\]
Assouad accumulates many binary testing difficulties across coordinates, while Fano uses an $M$-way classification problem.

The resulting lower bounds involve two competing quantities:
\[
\boxed{
\text{parameter separation}
\qquad \text{vs.} \qquad
\text{distributional closeness}.
}
\]
For example, Le Cam gives a lower bound schematically of the form
\[
R_n
\gtrsim
\|\theta_1-\theta_2\|^2
\left[
1-
\sqrt{
n\,\text{KL}(P_{\theta_1},P_{\theta_2})
}
\right].
\]

Thus, the selected parameters should be sufficiently far apart to force a large estimation loss, while their induced distributions should remain sufficiently close in total variation, KL divergence, or Hellinger distance so that they are difficult to distinguish statistically.

Finally, the scale of the construction, such as
\[
\Delta\theta,\qquad
\delta,\qquad
a,\qquad
m,\qquad
\epsilon,
\]
is chosen to maximize the resulting lower bound by balancing parameter separation against statistical distinguishability.

In summary,
\[
\boxed{
\begin{aligned}
&\text{Choose a hard finite subset of }\Theta
\\
&\qquad\Longrightarrow
\text{reduce estimation to testing}
\\
&\qquad\Longrightarrow
\text{bound testing error using TV/KL/Hellinger/entropy}
\\
&\qquad\Longrightarrow
\text{balance separation against distinguishability}
\\
&\qquad\Longrightarrow
\text{obtain a minimax lower bound}.
\end{aligned}
}
\]

If a matching estimator gives an upper bound of the same order, then that common order is the minimax rate.
\clearpage
\bibliography{ref}
\bibliographystyle{abbrv}
\end{document}
